\chapter{Introduction}
\label{ch:introduction}

\section{Motivation and objectives}
\label{sec:intro-problem}

This dissertation concerns inference for response-distribution features other
than the conditional mean: conditional quantiles and fixed-content intervals.
River-flow forecasts must describe lower and upper
tails as well as central behavior, and their accuracy can depend on the
forecasting system and lead time. Electricity-price forecasts must represent
asymmetric uncertainty conditional on recent prices and supplied covariates.
In a different setting, a tolerance interval must contain a prescribed fraction
of a population with a stated repeated-sampling confidence. These problems
involve related distributional ideas, but their inferential targets are not
interchangeable. The unifying question is how to define each target, construct
a model and computation appropriate to it, and validate the resulting claim
under the information actually available.

Conditional quantile regression provides a direct way to study tail and central
behavior. Its use in time series raises two recurring questions. First, how
should temporal dependence be represented when a linear state model is too
restrictive or when several forecast products carry different information?
Second, how can posterior computation, diagnostics, and forecasting remain
practical when the model must be fitted repeatedly or at several probability
levels? Fixed-content intervals raise a separate question: after probability
content has been specified, what criterion determines the placement of the
interval, and how should a data-dependent interval be calibrated?

This dissertation has four objectives:

\begin{enumerate}
\item develop and evaluate a dynamic quantile model that corrects
source-specific river-flow forecast products and combines separate quantile
fits into a predictive distribution;
\item provide statistical software and posterior computation for dynamic
quantile models and for static regularized regression under the
asymmetric Laplace (AL) and extended asymmetric Laplace (exAL) families,
including the high-dimensional readout fits used with recurrent features;
\item develop Bayesian conditional-quantile regression on fixed nonlinear
recurrent features, both separately and jointly across probability levels; and
\item characterize mean-tilted fixed-content intervals and study a calibrated
procedure for selecting a short tolerance interval, while identifying which
regression and dynamic extensions remain proposals.
\end{enumerate}

The resulting chapters address four collaborative research projects. Their
common subject is non-mean distributional inference, but each chapter has its
own model, uncertainty statement, and validation design. Together, the chapters
contribute statistical methods and computational results for conditional
quantiles and fixed-content intervals, while distinguishing model assumptions,
posterior approximations, predictive evidence, and repeated-sampling claims.

\section{Research contributions and antecedents}
\label{sec:intro-contributions}

\paragraph{Hydrologic forecast correction and synthesis.}
Chapter~\ref{ch:research-b} studies daily discharge for the San Lorenzo River
when gauge observations, retrospective hydrologic products, and issued
ensembles differ in support and forecast horizon. We develop a shared
dynamic quantile process with source-specific discrepancy states and combine
complete predictive distributions from separately fitted quantiles into a
nondecreasing marginal quantile function. We also develop joint
scale--skewness Laplace--Delta variational updates for the extended likelihood.
The exAL distribution, dynamic quantile linear models, forecast combination,
ensemble model output statistics, and monotone quantile correction retain
their cited origins. The empirical study uses five fixed forecast origins,
with 28-day and common eight-day comparisons. At each origin, forecasts
condition on one supplied precipitation and soil-water path. The evaluation is
restricted to five origins because each 28-day window requires seven
long-record dynamic quantile fits, extensive processing of multiple forecast
products, and substantial computational time, storage, and post-processing.
A matched full-model Markov chain Monte Carlo (MCMC) comparison remains a
priority for assessing variational approximation error and uncertainty.

\paragraph{Statistical computation and the \pkg{exdqlm} software.}
Chapter~\ref{ch:research-a} addresses the computation required by dynamic
quantile analysis. Yan, Zheng, and Kottas \citep{exd-yan2025family} introduced
the exAL family and static exAL regression. Barata, Prado, and Sans\'o
\citep{exd-me} developed the exDQLM formulation, its original MCMC and
importance-sampling variational Bayes (ISVB)
procedures, and the transfer-function extension. Our contribution is the
substantial later restructuring and extension of \pkg{exdqlm}:
improved Laplace--Delta variational Bayes (LDVB) and nonconjugate
scale--skewness updates; compiled implementations for intensive operations,
with \proglang{R} implementations retained where applicable; filtering and
smoothing corrections; forecasting, diagnostics, scoring, and
posterior-predictive synthesis; and static AL/exAL regression under
regularized-horseshoe shrinkage. The worked examples compare simulation cost,
approximation, uncertainty, and predictive loss under the reported models.

\paragraph{Bayesian quantile regression with deep echo-state features.}
Chapter~\ref{ch:research-c} develops the quantile deep echo state network
(Q--DESN) and its multiquantile extension (MQ--DESN). A deep echo-state
network maps response history and covariates into fixed recurrent nonlinear
features; these are conventionally called reservoir states. Bayesian
regression on those features estimates conditional quantiles without treating
the recurrent weights as unknown parameters. Q--DESN fits one probability
level at a time under AL or exAL likelihoods and shrinkage priors. MQ--DESN
uses a joint multiquantile generalized Bayes update with regularized-horseshoe
shrinkage on adjacent coefficient differences. This construction promotes
similar neighboring quantile functions; we do not impose explicit noncrossing
constraints. The evidence consists of
specified simulation mechanisms, one retrospective Global Flood Awareness
System (GloFAS) origin, and 114 electricity-price region--fold comparisons.
Reservoir specifications were selected using separate development series. The
simulation comparisons are conditional on those specifications and on one
evaluation trajectory per mechanism, rather than repeated simulation
replications.

\paragraph{Mean-tilted intervals and tolerance actions.}
Chapter~\ref{ch:research-d} begins from the fact that probability content does
not determine interval placement. Starting from the residual-product criterion
of \citet{rqr-PouplinEtAl2024RQR} and its population-content equation, the
project derives the retained-mean characterization of the mean-preserving
interval and embeds it in a retained-mean tilt family. It then develops a
profiled procedure that screens fitted content and tilt using a posterior
content probability before selecting the shortest eligible interval. Independent
simulation evaluates the prespecified procedure on eight named distributions
and characterizes its repeated-sampling performance over that panel. The
chapter also formulates regression and
dynamic extensions in which time-varying latent coefficients determine the
lower and upper endpoints. Simulation studies, numerical validation, and
applications of these extensions remain future work.

\section{Conditional quantiles and fixed-content intervals}
\label{sec:intro-targets}

Let \(Y\) be a response and let \(\mathcal F\) denote the information available
when a forecast is issued. For \(0<\tau<1\), the conditional \(\tau\)-quantile
is
\begin{equation}
Q_\tau(Y\mid\mathcal F)
=\inf\{q\in\mathbb R:\Pr(Y\le q\mid\mathcal F)\ge\tau\}.
\label{eq:intro-conditional-quantile}
\end{equation}
It minimizes conditional expected check loss. For \(u\in\mathbb R\),
\begin{equation}
\rho_\tau(u)=u\{\tau-\Ind(u<0)\},
\qquad
Q_\tau(Y\mid\mathcal F)
\in\argmin_{q\in\mathbb R}
\E\{\rho_\tau(Y-q)\mid\mathcal F\},
\label{eq:intro-check-loss}
\end{equation}
under standard regularity and uniqueness conditions
\citep{qdesn-KoenkerBassett1978,qdesn-Koenker2005}. The conditioning
information is part of the target: changing the response history, covariates,
or forecast origin can change the conditional quantile.

A model at one value of \(\tau\) estimates one conditional quantile. Chapters
\ref{ch:research-b} and~\ref{ch:research-a} fit several probability levels
separately and, when a full predictive distribution is needed, synthesize
their predictive output after fitting. Chapter~\ref{ch:research-c} studies
both separate Q--DESN regressions and MQ--DESN, a joint multiquantile
generalized Bayes construction that couples coefficients across levels.
Separate fits, post-fit synthesis, and a multiquantile generalized Bayes
posterior are therefore three distinct operations. Quantile crossing has been
addressed through rearrangement, isotonic projection, constrained estimation,
and monotone quantile-process models. Rearrangement and projection order a
fitted grid but do not combine separate parameter posteriors into a joint
posterior
\citep{qdesn-BarlowBrunk1972IsotonicDual,exd-chernozhukov2010,
qdesn-WuLiu2009StepwiseNCQR,qdesn-BondellReichWang2010NoncrossingQR,
qdesn-YangTokdar2017JointQuantilePlanes}.

A fixed-content interval has a different target. Let \(F\) be a continuous
population distribution and let \(c\in(0,1)\) be the required probability
content. Any endpoints \(a<b\) satisfying
\begin{equation}
F(b)-F(a)=c
\label{eq:intro-fixed-content}
\end{equation}
have the same content, but they need not have the same location or length.
For a right-skewed distribution, such as a lognormal distribution, an
equal-tailed interval allocates fixed tail probabilities on both sides, while
a shortest interval shifts toward the region of higher density. A
mean-preserving interval imposes yet another placement condition by requiring
the retained portion to preserve the population mean. Thus, content specifies
how much probability is retained, whereas a placement criterion determines
which interval is selected.

A tolerance statement concerns the repeated-sampling behavior of a
data-dependent interval. If \(\mathcal D\) denotes the random sample and
\([L(\mathcal D),U(\mathcal D)]\) is the reported interval, content \(c\) with
confidence \(1-\alpha\) requires
\begin{equation}
\Pr_{\mathcal D\sim F}
\left[
F\{U(\mathcal D)\}-F\{L(\mathcal D)^-\}\ge c
\right]\ge 1-\alpha.
\label{eq:intro-tolerance-statement}
\end{equation}
This is not implied by a credible interval for the endpoints or by a predictive
interval for one future response
\citep{rqr-Wilks1941ToleranceLimits,rqr-Wilks1942StatisticalPrediction,
rqr-Robbins1944DistributionFreeTolerance}. Interpretation depends on both the
interval construction and its validation criterion. A credible interval
assigns posterior
probability to a parameter, state, quantile location, or endpoint under a
specified update. A predictive interval concerns a future response under a
predictive distribution. A confidence interval is evaluated by repeated-sample
coverage of a parameter or functional. A tolerance interval is evaluated by
whether the complete data-dependent rule attains population content \(c\) with
confidence \(1-\alpha\).

\section{Models and posterior computation}
\label{sec:intro-updating}

The AL distribution is widely used in Bayesian quantile
regression because its location is a specified quantile and its negative log
density is proportional to check loss \citep{env-yu2001bayesian}. When the AL
distribution is used as a working likelihood rather than the true response
distribution, posterior concentration can still identify a quantile-regression
target, but ordinary posterior intervals need not have nominal frequentist
coverage
\citep{env-sriram2013theoretical,qdesn-YangWangHe2016ALInference,
qdesn-JiLeeRabeHesketh2025ValidSE}. The exAL family adds a shape parameter
while preserving the specified quantile and a conditionally Gaussian
representation. In their simulation study, \citet{exd-yan2025family} found
coefficient credible-interval coverage closer to nominal under exAL than AL,
most consistently at the 0.05 quantile.

Chapters~\ref{ch:research-b} and~\ref{ch:research-a} place the AL or exAL
quantile location in a Gaussian state-space model. Trend, seasonal, regression,
discrepancy, and transfer components can then evolve over time. Chapter
\ref{ch:research-c} instead constructs a fixed design from recurrent nonlinear
states and estimates the quantile-regression coefficients conditional on that
design. In both cases, latent-variable augmentation supports MCMC and
variational computation. Variational Bayes replaces parts of the posterior
integration by an optimized approximation and is useful for repeated or
long-series fits, but its tail and interval summaries remain approximate.

Chapter~\ref{ch:research-d} uses generalized Bayes rather than a response
likelihood for its interval endpoints. Let \(\theta\) denote the interval
parameter or action, \(\ell(\theta;y)\) the specified loss at observation
\(y\), \(\pi(\theta)\) a prior, \(y_{1:n}\) the observed sample, and \(w>0\) a
learning rate. The generalized posterior is
\begin{equation}
\pi_w(\theta\mid y_{1:n})
=
\frac{
\pi(\theta)\exp\{-w\sum_{i=1}^n\ell(\theta;y_i)\}
}{
\int \pi(\vartheta)
\exp\{-w\sum_{i=1}^n\ell(\vartheta;y_i)\}\,d\vartheta
}.
\label{eq:intro-generalized-bayes}
\end{equation}
The denominator normalizes the update when finite. This construction describes
uncertainty about a loss-defined interval action; it does not assert that
\(\exp\{-w\ell(\theta;y)\}\) is the sampling density of \(Y\)
\citep{rqr-BissiriHolmesWalker2016GeneralBayes}. Consequently, endpoint draws
are not posterior predictive response draws.

\section{Forecast information and evaluation}
\label{sec:intro-evaluation}

Forecasts must state which future inputs are available. In the hydrologic study
of Chapter~\ref{ch:research-b}, forecasts condition on one supplied
precipitation and soil-water path at each origin. The streamflow study in
Chapter~\ref{ch:research-c} conditions on
weather values observed retrospectively over the forecast window. Its
electricity-price study uses supplied day-ahead load, solar, and wind forecasts.
In the software
examples of Chapter~\ref{ch:research-a}, users must supply future design and
evolution matrices when forecasting beyond the stored model.

Evaluation criteria are matched to the estimated object. Check loss evaluates
one conditional quantile. The continuous ranked probability score (CRPS)
evaluates an entire predictive distribution \(G\) against an observation
\(y\):
\begin{equation}
\operatorname{CRPS}(G,y)
=
\E_{Y\sim G}|Y-y|
-\frac12\E_{Y,Y'\sim G}|Y-Y'|.
\label{eq:intro-crps}
\end{equation}
It is a proper, negatively oriented score, so smaller values indicate better
predictive performance; unlike check loss at one probability level, it assesses
the full predictive distribution \citep{exd-gneiting2007}.

Simulation can additionally assess recovery of a known parameter or quantile
path. Predictive-interval calibration, posterior convergence, and tolerance
content address different questions. The validation in
Chapter~\ref{ch:research-d}, for example, estimates the content and width of a procedure fixed
before the independent simulation stage. Across the prespecified
eight-distribution panel, it supplies finite-sample repeated-sampling evidence
for both content attainment and width. Evaluation under additional
distributions, dependence structures, and sample sizes requires separate
evaluation.
The application studies likewise define their comparisons through stated
forecast origins, information sets, and scoring rules.

\section{Dissertation organization}
\label{sec:intro-organization}

Chapter~\ref{ch:research-b} presents the hydrologic application and its
source-specific predictive synthesis; Chapter~\ref{ch:research-a} develops the
statistical computation and software used for dynamic and static quantile
models; Chapter~\ref{ch:research-c} develops Q--DESN and MQ--DESN; and
Chapter~\ref{ch:research-d} develops mean-tilted interval targets, a calibrated
tolerance procedure, and proposed regression and dynamic extensions.
Chapter~\ref{ch:synthesis} compares the four contributions, their limitations,
and the questions that remain. Selective appendices retain nonstandard
derivations, proofs, and claim-relevant supporting evidence.
